%%%%%%%%%%%%%%%%%%%%%%%%%%%%%%%%%%%%%%%%%
% Golang Reference Sheet
% Version 0.1 (01/09/2016)
%
% Author:
% Andy Walker (https://github.com/alaska)
% License:
% The MIT License (see included LICENSE file)
%
%%%%%%%%%%%%%%%%%%%%%%%%%%%%%%%%%%%%%%%%%

% Basic document setup
\documentclass[10pt]{article}
%\usepackage{lua-visual-debug}
\usepackage[portrait, margin=0.25in]{geometry}
\setlength{\parindent}{0pt}
\usepackage[T1]{fontenc}
\usepackage{fontspec}
\usepackage{luatexja-fontspec} %% Japanese text

% Local imports (and anything they include) first
\usepackage{import}
\subimport{theme/}{perldoc}

% Now that we have theme colors, we can set the main font
\setmainfont{Inconsolata}
\setmonofont{Inconsolata}
\setmainjfont{IPAexGothic}

% Package imports
\usepackage{tikz}
    \usetikzlibrary{tikzmark, calc, positioning, fit, shapes.misc}
\usepackage[outputdir=\directlua{tex.sprint(os.getenv('BUILD_DIR'))}]{minted}
\usepackage{tcolorbox}
\usepackage{pdfrender} %shadowed text
\usepackage{booktabs}

% Fix Minted spacing
\newlength{\fancyvrbtopsep}
\newlength{\fancyvrbpartopsep}
\makeatletter
\FV@AddToHook{\FV@ListParameterHook}{\topsep=\fancyvrbtopsep\partopsep=\fancyvrbpartopsep}
\makeatother
\setlength{\fancyvrbtopsep}{-3pt}
\setlength{\fancyvrbpartopsep}{0pt}

% Frame/box definitions
%\pagecolor{grPageColor}
%TODO: Look at sidebyside, overlay, tcbline, segmentation style
\newtcolorbox{codebox}[1][]{fontupper=\scriptsize\texttt, coltext=grText, colback=grBack, fonttitle=\scriptsize\texttt, coltitle=grHeadingText, colbacktitle=grHeadingBack, arc=3pt, boxrule=0pt, titlerule=0.5pt, box align=top, before skip=0pt, after skip=3pt, left skip=0pt, right skip=3pt, toptitle=3pt, bottomtitle=3pt, subtitle style={after skip=6pt}, boxsep=0pt, left=3pt, top=3pt, halign=left, title=#1}

% Code formatting alias/options
\usemintedstyle{perldoc}
\newminted{go}{autogobble}

% Custom commands
\newcommand*{\CheapShadow}[1]{%
  \begingroup
    \leavevmode
    \rlap{\kern.3pt\raise-.3pt\hbox{%
      \pdfrender{
        TextRenderingMode=FillStroke,
        LineWidth=.2pt,
        StrokeColor=black,
		FillColor=black,
      }#1%
    }}%
    \rlap{%
      \pdfrender{
        TextRenderingMode=Stroke,
        LineWidth=.2pt,
      }#1%
    }%
    \hbox{#1}%
  \endgroup
}

\begin{document}
\begin{minipage}[t][][b]{\textwidth}
    \fontfamily{pag}%
    \fontseries{b}%
    \textcolor{grTitle}{\CheapShadow{\large{Go Quick Reference}\tiny{version $\alpha$}}}\hfill
    \includegraphics[height=12pt]{./figures/gopher.eps}
\end{minipage}
\begin{minipage}[t]{.25\textwidth}
    \begin{codebox}[Basic Types]
    \tcbsubtitle{Slices}
    \tcbsubtitle{Maps}
    \tcbsubtitle{Channels}
    \tcbsubtitle{Type Declarations}
    \begin{gocode}
        //struct type
        type myType struct {
            embeddedType
            protectedMember int
            PublicMember    string
        }
        //type
        type intList []int
    \end{gocode}
    \tcbsubtitle{Functions}
    \begin{gocode}
        func returnNumberOne() int {
            return 1
        }
        // multiple return values
        func returnTwoInts() (int,int) {
            return 2,3
        }
        // methods have receievers
        func (t *myType) getProt() int {
            return t.protectedMember
        }
    \end{gocode}
    \tcbsubtitle{Interfaces}
    \begin{gocode}
        type protGetter interface {
            // function signatures only
            getProt() int
        }
    \end{gocode}
    \end{codebox}
    \begin{codebox}[Operators]
        \setlength{\tabcolsep}{2pt}
		\begin{tabular}{@{}l|llllllll@{}}
            Unary & +  & -  & ! & \textasciicircum & * & \& & <- &\tikzmark{a}\\
			Multiplication & * & / & \% & << & >> & \& & \&\textasciicircum&\\
			Addition & + & - & | & \textasciicircum &&&&\\
			Comparison & == & != & < & <= & > & >= &&\\
			Logical & \&\& &&&&&&&\\
            & || &&&&&&&\tikzmark{b}
		\end{tabular}
        \begin{tikzpicture}[remember picture, overlay]
            \draw [->,thick,draw=grAnnotate] ($(pic cs:a) + (0,5pt)$) -- ($(pic cs:b) + (0,-2pt)$) node [midway, sloped, above]{precedence};
        \end{tikzpicture}
        \textasciicircum x = bitwise compliment of x
    \end{codebox}
\end{minipage}%
\begin{minipage}[t]{.25\textwidth}
    \begin{codebox}[Declaration]
        \begin{gocode}
            var num int     // declare an integer
            var num int = 1 // declare & initialize
            var num = 1     // type inferred
            num := 1        // short declaration
            var num, foo int = 1, 2 // multiple
            // Unicode works
            var 名 = "アラスカ"
            // constant
            const name = "golang"
            // group-declaration
            var (
                foo int
                bar int = 1
                baz = 2
            )
            const (
                fooConst = iota // 0
                barConst        // 1
                bazConst = "Baz"
            )
            // camelCase when multi-word
            myFavoriteRodent = "gopher"
            // identifiers with uppercase first
            // letters are exported
            const PackageName = "my package"
            var BadIdea = "mutate from anywhere!"
        \end{gocode}
        \tcbsubtitle{Short Declaration}
        \begin{gocode}
            // Can only use in function body
            x := aFunction() // type inferred
            // can "re-declare" with at least one
            // new variable
            var y int
            x, z := 1, 2 // ok
            x, y := 3, 4 // ERROR
            // creates block-local variables
            for i:= 0 ; i < 10; i++ {
                fmt.Println(i)
            }
            fmt.Println(i) // ERROR
            // beware of shadowing
        \end{gocode}
        \tcbsubtitle{Iota}
        \begin{gocode}
        \end{gocode}
    \end{codebox}
\end{minipage}%
\begin{minipage}[t]{.25\textwidth}
	\begin{codebox}[Structure - Exectutable]
        \begin{gocode}
            package main
            import (
                "pkg"
                "pkg2"
            )
            import "fmt" //or one per-line
            const (
                numberOne = 1
            )
            const hello = "Hello, 世界"
            //main.main required for executable
            func main() {
                // Code
            }
		\end{gocode}
        \tcbsubtitle{Structure - Package}
        \begin{gocode}
            package mylib
            func CallMeFromOutside
        \end{gocode}
	\end{codebox}
    \begin{codebox}[Format verbs]
        Simpler than C's

        MOAR TABLE
    \end{codebox}
\end{minipage}%
\begin{minipage}[t]{.25\textwidth}
    \tcbset{right skip=0pt}%
    \begin{codebox}
        \begin{gocode}
            package anothermain
            import (
            "fmt"
            )
            func main() {
                fmt.Println("Gopher, save me")
            }
        \end{gocode}
    \end{codebox}
\end{minipage}
\end{document}
